% ---
% title: 擴散模型邊緣分布
% date: 2025/12/30
% tag: Diffusion
% excerpt: 用贝叶斯公式把当前时刻的 score 写成条件 score 的期望，以及这个等式在扩散模型里为什么好用。
% ---
当前时刻的 score function 可以通过条件分布来计算。

\[
\nabla \log p_t(x_t) = \mathbb{E}_{x_0 \sim p(x_0|x_t)}[\nabla \log p(x_t|x_0)]
\]

下面从贝叶斯公式和一个直观类比来看。

\section{贝叶斯视角}

贝叶斯公式给出： \[
p(x_t|x_0) p(x_0) = p(x_0|x_t) p(x_t)
\]

两边取对数，再对 \(x_t\) 求导： \[
\nabla_{x_t} \log p(x_t|x_0) + \nabla_{x_t} \log p(x_0) = \nabla_{x_t} \log p(x_0|x_t) + \nabla_{x_t} \log p(x_t)
\]

\(p(x_0)\) 不依赖 \(x_t\)，第二项为 0。移项得： \[
\nabla_{x_t} \log p(x_t) = \nabla_{x_t} \log p(x_t|x_0) - \nabla_{x_t} \log p(x_0|x_t)
\]

对 \(x_0\) 在 \(p(x_0|x_t)\) 下取期望： \[
\nabla \log p_t(x_t) = \mathbb{E}_{x_0 \sim p(x_0|x_t)}[\nabla \log p(x_t|x_0)]
\]

\section{直觉理解}

边缘分布的 score 是条件分布 score 的期望。左边 \(\nabla \log p_t(x_t)\) 是在所有可能的干净数据 \(x_0\) 下，\(x_t\) 的总体梯度方向。右边对每个可能的 \(x_0\) 算 \(\nabla \log p(x_t|x_0)\)（前向过程的 score），再按后验 \(p(x_0|x_t)\) 加权平均。

一个直观的类比：把 \(x_t\) 想成一个观测到的噪声数据点。它可能来自多个不同的干净数据 \(x_0\)，每个 \(x_0\) 都会把 \(x_t\) 往某个方向拉，这就是 \(\nabla \log p(x_t|x_0)\)。最后的 score 是这些拉力的加权平均，权重是「这个 \(x_0\) 真的是源头」的概率。

\section{扩散模型中的应用}

扩散模型里 \(p(x_t|x_0)\) 通常是高斯，所以： \[
\nabla \log p(x_t|x_0) = -\frac{x_t - \mu_t(x_0)}{\sigma_t^2}
\]

于是： \[
\nabla \log p_t(x_t) = \mathbb{E}_{x_0}[-\frac{x_t - \mu_t(x_0)}{\sigma_t^2}] = -\frac{x_t - \mathbb{E}[\mu_t(x_0)]}{\sigma_t^2}
\]

所以 score function 实际上指向条件均值的期望方向。

这也是训练去噪模型 \(\epsilon_\theta\) 等价于学 score function 的原因。
